\cotefile{57}
\scanpage{164}

\[
Y\subset\overset{\substack{\unsure{x}\\\tikz[baseline=-.4ex]\node[inner sep=0pt,rotate=-90] {$\scriptstyle\in$};}}{X}
\]

\diagram{\begin{tikzcd}[ampersand replacement=\&,column sep=small,row sep=small]
Y\subset X\arrow[dd,dash] \& X\times\mathrm{Pic}\arrow[d] \& {}\\
{} \& \mathrm{Pic} \& \mathrm{Div}\arrow[l]\\
\cdot \& p \& d
\end{tikzcd}}

\[
\underline L^{Y}
\]

\[
\xi\in\underline H^{0}_{f}(\underline L\unsure{Y|^{Y}})\qquad\text{\unsure{non nul}}
\]

\[
0\unsure{\longrightarrow|\mathord{-}}\mathrm{Ker}\longrightarrow\underline H^{0}_{f}(\underline L Y)^{\vee}\longrightarrow\underline O_Y\longrightarrow 0
\]

\[
0\longrightarrow\underline O_S\longrightarrow\underline H^{0}_{f}(\underline L)\longrightarrow\underline H^{0}_{f}(\underline L\mid Y)
\]

\[
\underline J\underline L\longrightarrow\underline L\longrightarrow\underline L\mid Y\longrightarrow0
\qquad\qquad \underline L=\underline J^{-1}
\]

\[
0\longrightarrow\underline O_X\longrightarrow\underline L\longrightarrow\underline L\mid Y\longrightarrow0
\]

\[
\begin{array}{ccccccccc}
0&\longrightarrow&\underline O_S&\longrightarrow&\underline H^0_f(\underline L)&\longrightarrow&\underline H^0_f(\underline L\mid Y)&\longrightarrow&\underline H^1_f(\underline O_X)\\
&&&&\text{tangent relatif}&&\text{Tangent}&&\text{Tangent}\\
&&&&\text{\unsure{de} }\mathrm{Div}/\mathrm{Pic}&&\mathrm{Div}/S&&\mathrm{Pic}/S
\end{array}
\]

\[
\begin{array}{rclcl}
\longrightarrow&\underline H^1_f(\underline L)&\longrightarrow&\underline H^1_f(\underline L\mid Y)&\longrightarrow\underline H^{1}_f\unsure{(\underline O_X)|\underline O_X}\\
&\text{son annulation universelle}&&\text{son annulation}&\\
&\text{implique que }\mathrm{Div}/\mathrm{Pic}&&\text{universelle}&\\
&\text{est simple le long de}&&\text{implique}&\\
&\text{la section \unsure{canonique}}&&\text{que }\mathrm{Div}&\\
&&&\text{est simple}&\\
&&&\text{sur }S\text{ le long}&\\
&&&\text{de la \unsure{section}}&
\end{array}
\]

% archive sheet number: 163

\scanpage{165}

\mspage{- 34 -}\hfill n$^{\mathrm{o}}$ 297

Soient $A$ un anneau topologique séparé, $E$ un $A$-module topologique séparé. Soit $\hat E$ le groupe additif complété du groupe topologique commutatif $E$ (\S\,3, n$^{\mathrm{o}}$5, th.2). L'application $\handwritten{\unsure{Z|\mathbb Z}}$-bilinéaire $(\lambda,x)\longrightarrow\lambda x$ du produit $A\times E$ des groupes additifs $A,E$ dans le groupe additif $E$ se prolonge par continuité en une application $\handwritten{Z}$-bilinéaire de $\hat A\times\hat E$ dans $\hat E$ (n$^{\mathrm{o}}$5, th.1), que nous désignerons encore par $(\lambda,x)\longrightarrow\lambda x$. En vertu du principe de prolongement des identités, on a encore $\lambda(\mu x)={}=(\lambda\mu)x$ pour $\lambda\in\hat A$, $\mu\in\hat A$, $x\in\hat E$ ; la loi externe $(\lambda,x)\longrightarrow{}\longrightarrow\lambda x$ définit donc sur $\hat E$ une structure de $\hat A$-module compatible avec la topologie de $\hat E$ ; si en outre $A$ admet un élément unité $1$, cet élément est aussi élément unité de $\hat A$, et si $E$ est un $A$-module unitaire, on a encore $1.x=x$ pour $x\in\hat E$ par le principe de prolongement des identités. Nous dirons que le $\hat A$-module topologique $\hat E$ ainsi défini est le \uline{complété} du $A$-module topologique $E$\unsure{,|.}

Soit $E$ un $A$-module topologique, $A$ et $E$ n'étant pas nécessairement séparés ; soient $N$ (resp. $F$) l'adhérence de $\{0\}$ dans $A$ (resp. $E$) ; $N$ est un idéal bilatère de $A$ (n$^{\mathrm{o}}$3; prop.6) et $F$ un sous-$A$-module de $E$ (n$^{\mathrm{o}}$1, prop.1) ; en outre on a par continuité $\lambda x\in F$ lorsque $\lambda\in N$ ou $x\in F$. On en déduit aussitôt, par passage aux quotients une application $(\lambda^{\circ},\overset{\circ}{x})\longrightarrow\lambda^{\circ}\overset{\circ}{x}$ de $(A/N)\times(E/F)$ dans $E/F$, et on vérifie aussitôt (grâce à la prop.24 du \S\,2, n$^{\mathrm{o}}$9) que cette application est loi de composition externe d'une structure de $(A/N$-module topologique sur $E/F$. Si l'on pose $B=A/N$, $L=E/F$, on dit que le $B$-module $L$ est le module séparé \uline{associé} à $E$; son complété $\hat L$ est un $\hat B$-module topologique, qu'on appelle le module \uline{complété séparé} de $E$.

% archive sheet number: 164

\scanpage{166}

AK

\uline{Pseudomorphismes}

$X$ un préschéma. ouvert schématiquement dense $U$ : si $X$ est le plus \deleted{petit} ss-préschéma fermé de $X$ dominant $U$ i.e. si $\overline U=X$ [lorsque $U\longrightarrow X$ est quasi-cpt cela implique que pour tout ouvert $V$, $U\cap V$ est schématiquement dense dans $X$ -- aussi cela vaut-il mieux comme définition : toute \deleted{fonction}\inserted{section de $\underline O_X$} sur un ouvert $V$ de $X$ nulle sur $U\cap V$ est nulle]

% slanted
\margin{\unsure{condition équiv. à ci-dessous}}

\deleted{\unsure{L'intéresse|L'intérêt}} Caractère local de la notion [Si $X=\bigcup X_i$, les $X_i$ ouverts, alors $U$, schématiquement dense ds $X\Longleftrightarrow$ les $U\cap X_i$ sch. dense dans les $X_i$]

$U$ schémati. dense dans $X\Rightarrow U$ schém. dense dans \deleted{des} $V$ \unsure{q.c.q.|q.p.|q.c.} $U\subset V\subset X$

Si $U,U'$ sont schém. denses, leur $U\cap U'$ l'est

[Soit $X$ et $Y$ sur $S$, $\deleted{X}Y$ séparé sur $S$, $\deleted{X}\unsure{\Rightarrow|\longrightarrow}U$ schém. dense dans $X$, si $f,g:X\longrightarrow Y$ sont des $S$-morphismes tels que $f\mid U=g\mid U$

\deleted{\unsure{Quasi}}\inserted{Pseudo}-morphismes : \unsure{classes|classe} d'équivalence d'\deleted{applications}morphismes définis sur des ouverts schématiquement denses, comme $S$-morphismes. Pour un quasi-$S$-morphisme dans un $S$-\deleted{pré}\uline{schéma}, il y a un \inserted{\uline{(plus grand)}} domaine de définition.

\[
\deleted{\text{Le domaine }D(f\mid V)={}}\quad
\mathrm{Dom}(f\mid V)=\mathrm{Dom}(f)\cap V
\]

% archive sheet number: 165

\scanpage{167}

\uline{Ex} Si $X$ \deleted{\unsure{est} \illegible{1} \unsure{contient} \unsure{le} \illegible{2} irréductibles} \inserted{est \uline{réduit}, alors schématiquement dense = dense}, pseudo-morphismes $=$ applications rationnelles

Pour tout ouvert $U$ de $X$, soit $F(U)$ l'ens des pseudo-$S$-morphismes de $U$ dans $Y$ ($S$-schém. séparé). Application \deleted{de}restriction $F(U)\unsure{\mathord{-}|\longrightarrow}F(\unsure{V|U})$ si $U\supset V$. Cela \deleted{est} définit un \uline{faisceau} [grâce à la relation $\mathrm{Dom}(f_i\mid U_{ij})=\mathrm{Dom}(f_i)\cap U_{ij}$], appelé \uline{faisceau des} \deleted{\unsure{quasi-morph.}} \uline{pseudo-$S$-morphismes} de $X$ dans $Y$.

En particulier, on obtient ainsi le faisceau des pseudo-sections de $\underline O_X$ (\deleted{\unsure{f}}faisant $X=S$, $Y=X[t]$) qui est un faisceau d'anneaux \uline{\unsure{contenant|$\longleftarrow$}} $\underline O_X$, en particulier un faisceau d'algèbres.

\deleted{\uline{Prop.} 0.1 Ce faisceau est quasi-cohérent. On peut supposer $X=\mathrm{Spec}(A)$.}

% archive sheet number: 166

\scanpage{169}

\margin{Diviseur de fonctions}

Soit $X$ une variété irréductible \uline{normale} déf/$k$ $K=k(X)$ le corps des fonctions rationnelles/$k$ sur $X$

Soit $V$ une \deleted{\unsure{hypersurface irréductible}}\inserted{diviseur premier déf.} sur $k$ sur $X$

Alors l'ensemble des $f\in K$ \uline{définies} en $V$ forment un anneau de valuation discrète, $\deleted{\unsure{\underline O/V}}\underline O_{(V)}$.

En effet, c'est une \deleted{\illegible{1}} localité normale, dont la dimension de Krull est $1$.

Soit $\unsure{v|\nu}_V$ la valuation \inserted{de $K$} attachée à cet anneau de valuation du $K$. \unsure{Appelont|Appelons} \uline{diviseur} \inserted{déf/$k$} de $X$ toute comb. lin. formelle des diviseurs premiers \unsure{sur|rat.} déf/$k$. Posons, pour toute $f\in K$

\[
\mathrm{div}\,f=\sum_{\substack{V\ \mathrm{div.}\\ \mathrm{pr}/k}}v_V(f).\,V
\]

Cette somme a un sens, car pour $f$ donnée, il n'y a qu'un nb fini de termes non nuls ; \deleted{car \unsure{tout}} pour le voir, on est ramené au cas \deleted{\unsure{ouvert}}, affine, \unsure{alors|tous} les diviseurs pr. déf/$k$ correspondent aux idéaux premiers non nuls minimaux de l'anneau normal noethérien $A=k[X]$, et la chose est bien connue.

\uline{Théorème} \deleted{\unsure{pour} \illegible{1} \unsure{diviseur principal|diviseur équivalent}} \inserted{$\mathrm{div}\, f\geq 0$, $\Longleftrightarrow f$ est partout définie}

\unsure{Ceci|Cor.|C'est} résulte de la théorie des anneaux normaux : $A=\bigcap A_{\mathfrak p}$ ($\mathfrak p$ premier minimal) si $A$ est normal. En effet, on est aussitôt ramené au cas affine, et $\mathrm{div}\,f\geq0\Longleftrightarrow f\in\bigcap A_{\mathfrak p}$.

% archive sheet number: 168

\scanpage{170}

\uline{Corollaire 1} \deleted{Pour que $f$ soit définie en une \unsure{la} sous-variété irréductible $Y$, il faut et il suffit que $Y$ ne soit \unsure{contenu|contenue} dans aucun \unsure{des} \unsure{$V$} tels que} \inserted{L'ens. des pts de $X$ où $f$ n'est pas définie est la réunion des supports des $V$ tels que $v_V(f)<0$.} C'est \unsure{encore vrai} si $f$ est une application rationnelle de $X$ dans une variété \deleted{\uline{affine}} ($X$ normal) c'est par suite une sous-variété pure de dim $n-1$.

\deleted{\unsure{Th.}} N.B. Les résultats précédents valent sans changement essentiel si \deleted{\unsure{X n'est pas}} \inserted{l'ens} des pts singuliers de $X$ est de dim $\leq n-2$.

\uline{Th.} Soit $X$ une variété \deleted{non singulière}\inserted{normale de dim $n$}, $Y$ \deleted{\unsure{q}} une variété complète, $f:X\longrightarrow Y$ une application rationnelle. Alors l'ens. des pts où $f$ n'est pas définie est de dim $\leq n-2$.

\deleted{On peut supposer $X$} Deux démonstrations

a) On peut supposer, en regardant le \uline{graphe} de $f$, que $f$ est inverse d'un \uline{morphisme} birationnel $g:Y\longrightarrow X$, et $g$ est \deleted{propre} \uline{propre}.

\deleted{L'ens} Soit $X'$ l'ens des pts $x$ \deleted{\unsure{de}} tels que $g^{-1}(x)$ soit de dim $>0$, alors $X'$ est de dim $\leq n-2$, car sinon $g^{-1}(X')$ serait de dim $\geq(n-1)+1=n$, absurde. \deleted{Donc $g^{-1}(X)$} D'autre part, il résulte du Main th. qu'au dessus de $\complement X'$, $g$ est un isom. birégulier, donc $g^{-1}$ est définie sur $\complement X'$. O.K.

N.B. Bien entendu, le th. se généralise aussi au cas où \deleted{\unsure{$X$}} l'ens des singularités de $X$ est de dim $\leq n-2$

% archive sheet number: 169

\scanpage{171}

On suppose $X$ irréductible, $f$ \unsure{dominant|dominante|domine}, ce qui est loisible.

b) Soit $V$ une hypersurface irréductible de $X$. Pour voir que $f$ est défini en au moins un pt de $V$, \deleted{\unsure{il su}} i.e. $f$ ``défini en la localité $\underline O_{(V)}$'', il faut montrer qu'il existe une localité de $Y$ dominée par $\underline O_{(V)}$. Or cela résulte de la définition valuative des variétés complètes, et du fait que $\underline O_{(V)}$ est un anneau de valuation.

% archive sheet number: 170

\scanpage{173}

\uline{Le foncteur de Picard}

\margin{Faits élémentaires sur les classes de diviseurs, faisceaux inversibles, leurs propriétés etc}

Soit $X$ un $S$-préschéma, on pose

\[
\mathrm{Prépic}\,(X/S)=\deleted{H^1(X,\underline O_X^*)}\mathrm{Pic}(X\deleted{\unsure{/S}})=H^1(X,\underline O_X^*)
\]

\[
\text{\uuline{Prépic}}_{X/S}(S')=\mathrm{Prépic}\,(X'/S')=H^1(X',\underline O_{X'}^*)\qquad\text{où }\ X'=X\times_S S'
\]

donc $\text{\uuline{Prépic}}_{X/S}$ est bien un foncteur \uuline{contravariant} en le préschéma $S'$

\[
\mathrm{Prépic}'(X/S)=\Gamma(S,\underline H^1_{X/S}(\underline O_X^*))
\]

\[
\text{\uuline{Prépic}}'_{X/S}(S')=\Gamma(S',\underline H^1_{X'/S'}(\underline O_{X'}^*))\qquad\text{où }\ X'=X\times_S S'
\]

% slanted
\margin{expliciter les objets, l'\unsure{équivalence}, \unsure{pour} $\mathrm{Pic}(X\deleted{/S})$ définir \unsure{fn} $\mathrm{Pic}(X\deleted{/S})\longrightarrow\mathrm{Prépic}'(X/S)$}

On passe de \uuline{Prépic} à \uuline{Prépic}$'$ en exigeant la \unsure{descente} par recouvrements ouverts.

\uline{Proposition 1} Supposons $\underline H^0_{X/S}(\underline O_X)=\underline O_S$. Alors on a un \deleted{isomorphisme}\inserted{\uline{monomorphisme}} canonique

% slanted
\margin{mettre \deleted{\unsure{ce} corollaire} sous forme d'une suite exacte plutôt}

\[
\mathrm{Prépic}'(X/S)\mathrel{\deleted{\unsure{\sim}}\!\longleftarrow}\mathrm{Prépic}\,(X/S)/\mathrm{Pic\,abs}\,(S)
\]

qui est un isomorphisme si $X/S$ a une section Résulte de l'assertion plus précise

\uline{Corollaire 1} Le foncteur $\underline F\longrightarrow\underline F\otimes_S X$ des faisceaux loc. libres de $S$ dans les faisceaux loc. libres de $X$ \deleted{définit un}\inserted{est} \uline{pleinement} fidèle, l'image essentielle est formée des faisceaux loc. libres sur $X$ qui sont triviaux loc. au dessus de $Y$, le foncteur inverse sur ce dernier est le foncteur $\underline H^0_{X/\unsure{S|B}}$.

% slanted
\margin{\unsure{cela} pour un \unsure{cert.|cond:} \unsure{groupe|graphe} ? ou une proposition préliminaire}

% slanted
\margin{dire que les \unsure{images} sont \unsure{plus} \unsure{petits gr.}}

% archive sheet number: 172

\scanpage{174}

Il en résulte en effet que l'on a une suite exacte (en se bornant aux faisceaux inversibles)

\[
0\longrightarrow\mathrm{Pic}(S)\longrightarrow\mathrm{Pic}(X)\longrightarrow\mathrm{Prépic}'(X/S)
\]

% slanted
\margin{[surjectif si $X/S$ a une section] les \unsure{hom} sont définis en tous cas, mais l'exactitude demande que $\underline H^0_{\unsure{X/S|f/S}}(\underline O_X)=\underline O_S$}

l'\inserted{\unsure{dernière|dernier}}image étant formée des sections de $\underline H^1_{X/S}(\underline O_X^*)$ qui sont dans l'image de $\underline H^1(\underline O_X^*)$ i.e. dont l'obstruction dans $H^2(S,\underline O_S^*)$ est nulle

[N.B. le début de la suite exacte n'est que la suite exacte de \unsure{bas degré} de la suite spectrale de Leray

\diagram{\begin{tikzcd}[ampersand replacement=\&,column sep=tiny,row sep=small]
0\arrow[r]\&H^1(S,\underline H^0(\underline O_{X\deleted{\unsure{*}}}^*))\arrow[r]\&H^1(X,\underline O_{X\deleted{*}}^*)\arrow[r]\&H^0(S,\underline H^1(\underline O_{X\deleted{*}}^*))\arrow[r]\&H^2(S,\underline H^0(\underline O_X^*))\arrow[d]\\
{}\&{}\&{}\&{}\&H^2(X,\underline O_{X\unsure{*}}^{\unsure{*}})
\end{tikzcd}}

compte tenu que $\underline H^0(\underline O_X^*)=\underline O_S^*$]. S'il y a une section \inserted{$s$} de $X/S$, alors on \deleted{aura}\inserted{a} \deleted{un hom inverse} un résultat plus complet

\deleted{$\unsure{s^*}$.}

\uline{Corollaire 2} Soit $\underline C$ la catégorie des faisceaux inversibles $\inserted{\underline L}$ sur $X$ munis d'un \uline{isomorphisme} $\underline O_S\longrightarrow s^*(\underline L)$. Alors \deleted{\unsure{pour}}\inserted{(i)} tout automorphisme d'un objet de $\underline C$ est trivial (ii) Tout \deleted{tout} élément de $\mathrm{Prépic}\,(X/S)$ provient d'un $\underline L\in\underline C$, déterminé à isom. \uline{unique} près.

% archive sheet number: 173

\scanpage{175}

Pour prouver (ii), on considère un recouvrement de $S$ par des $U_i$ tels que $\unsure{\zeta|\xi}$ soit donné par des $\underline L_i$ sur $X_i=X\mid U_i$ ; prenant les $U_i$ assez petits

\deleted{[envisageant p.ex. $\unsure{s\in S|x\in S}$] on peut supposer que}

Remplaçant $\underline L_i$ par $\underline L'_i\otimes_S s^*(\underline L_i)^{-1}$, on peut supposer donnés des isom $\underline O_S\mid U_i\xrightarrow{\overset{\varphi_i}{\sim}}s^*(\underline L_i)$. Dans $X_{ij}=X\mid U_i\cap U_j$, $\underline L_i$ et $\underline L_j$ sont \uuline{$U_{ij}$-équivalents},

% slanted
\margin{terminologie à introduire}

donc il y a un isomorphisme \uline{unique} $\varphi_{ij}$ de $\underline L_i$ sur $\underline L_j$ compatible avec les $\varphi_i$. \uline{De} l'unicité résulte que les $\varphi_{ij}$ satisfont la condition de recollement, d'où le $\underline L$ cherché. Question : donner une formulation cohomologique des raisonnements.

% slanted
\margin{en tous cas, il faut harmoniser avec la démonstrat\textsuperscript{ion} de \unsure{prop 2|cor 2} (langage de la descente)}

\deleted{\uline{\unsure{Corollaire 2}}}\inserted{\uline{Proposition 2}} On suppose encore $\underline H^0_{X/S}(\underline O_X)=\underline O_S$, et $X$ quasi-compact séparé sur $S$, \unsure{donc} \deleted{est} relation précédente satisfaite aussi après extension plate de la base. Soit $S'\longrightarrow S$ une extension fid\textsuperscript{t} plate \inserted{qu.cpte} de la base. i) Alors

\[
\deleted{\unsure{\text{est}}}\qquad\deleted{\mathrm{Pic}}\mathrm{Prépic}'(X/S)\longrightarrow\mathrm{Prépic}'(X'/S')
\]

est injectif. (ii) Si $X/S$ a \uline{localmt} \uline{sur $S$} une section, alors la suite d'ensembles

% archive sheet number: 174

\scanpage{176}

\diagram{\begin{tikzcd}[ampersand replacement=\&,column sep=small]
\deleted{\unsure{\text{on}|\longrightarrow}}\ \mathrm{Prépic}'(X/S)\arrow[r]\&\mathrm{Prépic}'(X'/S')\arrow[r,shift left=.65ex]\arrow[r,shift right=.65ex]\&\mathrm{Prépic}'(X''/S'')
\end{tikzcd}}

[où $S''=S'\times_S S'$] est exacte, \deleted{i.e. la suite d'ens $\text{\uuline{Prépic}}_{X/S}$}

\deleted{\unsure{soit}} En effet, le point (i) est encore une conséquence immédiate de

\uline{Corollaire 1} Sous les conditions de prop. 1, \inserted{$X$ séparé quasi-cpt sur $S$} corollaire 1, soit $S'\longrightarrow S$ une extension fid plate \inserted{et qu. cpte} de la base. Soit $\underline E$ \deleted{\unsure{f}}faisceau \deleted{inversible}\inserted{loc. libre} sur $X$.

Pour qu'il soit isomorphe à l'image inverse d'un faisceau loc libre \deleted{[\unsure{rappelons que}]} sur $S$, il \unsure{f et s|faut} que $\underline E'$ soit isomorphe à l'image inverse d'un faisceau loc libre sur $S$.

En effet, il faut exprimer que $f_*(\deleted{\unsure{F}}\underline E)$ est loc. libre et $\deleted{\unsure{g}}f^*f_*(\underline E)\longrightarrow\underline E$ un isom. Or $f_*(\underline E)'=f_*(\underline E')$ \inserted{(cf \unsure{Künneth} \unsure{plat|plate})} et on est ramené à des assertions immédiates de stabilité de propriétés par extensions fid. plates (\unsure{Exp|Chap.} \uline{IV})

% archive sheet number: 175

\scanpage{177}

Pour prouver (ii), on est ramené à prouver que t\textsuperscript{t} élément de $\mathrm{Prépic}'(X'/S')$ dont les images sont les \unsure{mêmes|$\text{m}^{\mathrm{es}}$}, provient d'un élément de $\mathrm{Prépic}'(X/S)$. À cause de l'unicité, on voit que c'est local sur $S$, on peut alors supposer qu'on a une section de $X$ sur $S$.

On peut alors identifier les \deleted{objets}\inserted{élémts} de $\mathrm{Prépic}\,(X/S)$ à des faisceaux inversibles \inserted{(déts à \uline{isom unique} près)} sur $X$ à section marquée sur $S$, ce qui permet d'appliquer la théorie de la descente \uline{fid plate}.

\uline{Corollaire 2} \unsure{Bien|On} \unsure{combiner|combine} avec Prop 1. (ii) pour donner des \uline{faisceaux} sur $X$.

\uline{Corollaire \uline{3}} Supposons que la relation $H^0_{X/S}(\underline O_X)=\underline O_S$ reste vraie pour tte extension de la base (i.e., $X$ \unsure{coh.|est} plat sur $S$ en dim 0, et $\underline H^0_{X/S}(\underline O_X)=\underline O_S$), \inserted{et $X/S$ a localement une section} Alors le foncteur $\text{\uuline{Prépic}}'_{X/S}$ est compatible avec [= exact pour] \deleted{\unsure{une}} la descente fidèlement plate quasi-compacte.

% archive sheet number: 176

\scanpage{178}

\deleted{\uline{\unsure{Prop.}} \unsure{Cor. 3}}

\[
\begin{aligned}
\text{\uuline{Pic}}_{X/S}&=\text{enveloppe exacte de }\text{\uuline{Prépic}}_{X/S}\ (\text{ou }\text{\uuline{Prépic}}'_{X/S})\\
&\text{pour la descente fid plate quasi-compacte}
\end{aligned}
\]

[Il y a un ennui, car un \uline{recouvrement} ouvert d'un préschéma non qu. cpt. n'est en général pas majoré par un morphisme fid plat qu. cpt. Il faut \unsure{donc|sans} dire au lieu de fid plat qu. cpt : \deleted{fid plat après localisation} \uline{morphismes de localisation fidèlement plats}]

Donc on a des morphismes fonctoriels

\diagram{\begin{tikzcd}[ampersand replacement=\&]
\text{\uuline{Pré\deleted{\unsure{P}}pic}}_{X/S}\arrow[r,"\alpha"]\&\text{\uuline{Prépic}}'_{X/S}\arrow[r,"\beta"]\&\text{\uuline{Pic}}_{X/S}
\end{tikzcd}}

Le premier morphisme n'est pas injectif en général [son noyau, si $\underline H^0_{X/S}(\underline O_X)=\underline O_S$, et si cela reste vrai par changmt de base] est le foncteur \uuline{Pic} (absolu) en vertu de prop. 1, (i), \inserted{[qui est nul \unsure{en passant sur un schéma|sur les schémas} \unsure{local|locaux}]} Il est surjectif si on suppose en plus de la condition de [~] que $X/S$ a une section. \deleted{\unsure{Pour avoir la}}

% archive sheet number: 177

\scanpage{179}

D'autre part $\beta$ est un mon. sous \unsure{la s\textsuperscript{l}|la seule} hypothèse $\underline H^0_{X/S}(\underline O_X)=\underline O_S$ ``universelle'' ; et il est surjectif donc bijectif si on suppose \uline{de plus} que \deleted{\unsure{$\mathrm{Pic}$}} $X$ a localmt une section sur $S$.

\uline{Exemple} où $\beta$ n'est pas surjectif : variété de \unsure{Sévéri|Sé véri|Severi}-Brauer non triviale sur un corps, \ldots\ On verra qu'il y a une \inserted{\unsure{théorie}} d'obstruction à la surjectivité de $\beta\alpha$, à valeurs dans un $H^2(S,G_m)$ \ldots.

\uline{Remarques}

On voit donc que le foncteur \uuline{Prépic} n'est pas représentable \deleted{en général, que \unsure{dans un}}\inserted{sauf dans des cas} \deleted{\unsure{pas|par}} triviaux sans intérêt ; $\mathrm{Prépic}'_{X/S}$ ne peut être représentable que si $\beta$ est un isom et si $\text{\uuline{Pic}}_{X/S}$ est représentable, \deleted{\unsure{(que} \unsure{si} $\unsure{g}$} [ce qui en pratique joint à des conditions comme celles déjà envisagée, y inclus l'existence d'une section de $X$ sur $S$].

Quant à $\text{\uuline{Pic}}_{X/S}$, il est représentable dans des cas \unsure{généraux|suffisamment} importants qui seront déterminés par la suite. C'est par suite lui le plus important pour notre propos.

% archive sheet number: 178
